% Documento fittizio per provare la navigazione fra pagine.
%
% Ogni pagina porta un numero enorme al centro: serve a capire a colpo d'occhio
% se le frecce ti hanno portato dove pensavi, senza dover leggere il testo.
% Le pagine hanno contenuti diversi apposta - testo fitto, una tabella, una
% figura geometrica - così si vede anche come rende la rasterizzazione.
\documentclass[12pt,a4paper]{article}
\usepackage[utf8]{inputenc}
\usepackage{lmodern}
\usepackage[T1]{fontenc}
\usepackage[italian]{babel}
\usepackage{geometry}
\geometry{margin=2.5cm}
\usepackage{xcolor}
\usepackage{tikz}

% Il numerone in cima alla pagina. Contenuto normale e non un overlay tikz: un
% overlay come primissima cosa del documento si porta via una pagina tutta sua,
% e il numero finiva sfasato rispetto al contenuto.
\newcommand{\numerone}[1]{%
  {\centering\fontsize{90}{90}\selectfont\bfseries\textcolor{gray!30}{#1}\par}%
  \vspace{0.5em}%
}

\title{Manuale di prova}
\author{Cartella dimostrativa di CleeCode}
\date{}

\begin{document}

% ---- 1 -------------------------------------------------------------------
\maketitle
\thispagestyle{empty}
\numerone{1}
\begin{center}
\Large Se stai leggendo questa pagina, la prima si è aperta da sola.\\[1em]
\large Premi $\rightarrow$ per andare avanti, $\leftarrow$ per tornare.
\end{center}
\vfill
\noindent In basso a destra nel riquadro dovresti vedere \texttt{pagina 1 di 6}.
Se dice solo \texttt{pagina 1}, vuol dire che né \texttt{pdfinfo} né
\texttt{ghostscript} hanno saputo dire quante pagine ci sono: sfogliare funziona
lo stesso, si ferma da solo quando le pagine finiscono.
\newpage

% ---- 2 -------------------------------------------------------------------
\numerone{2}
\section{Testo fitto}
Questa pagina serve a vedere se il corpo del testo resta leggibile alla
risoluzione con cui viene rasterizzata la pagina.
Nel mezzo del riquadro di un editor che gira dentro un terminale, la domanda
non è se una pagina si possa disegnare, ma chi la disegna. Un programma che
scrive dentro un riquadro non parla mai col terminale ospite: parla con
l'emulatore che sta in mezzo, e quello traduce tutto in una griglia di celle
prima di ridipingerla. Le sequenze grafiche non sopravvivono a quel passaggio,
perché il parser non le riconosce affatto e le lascia cadere.

Da qui, invece, i byte escono sullo stesso canale su cui esce tutto il resto
dell'interfaccia, e arrivano al terminale ospite intatti. È una differenza di
percorso, non di qualità del terminale, ed è anche il motivo per cui la stessa
immagine continua a vedersi attraverso una sessione remota: quello che viaggia
sul filo è l'uscita dell'editor, non quella di un programma dentro una shell.

La risoluzione con cui questa pagina è stata rasterizzata è la stessa per tutte
le pagine del documento, e viene poi ridimensionata per stare nel riquadro. Se
il testo qui sotto si legge senza fatica, il compromesso è quello giusto; se
sfarfalla, il numero da cambiare è uno solo, e sta accanto alla ragione per cui
è stato scelto.
\newpage

% ---- 3 -------------------------------------------------------------------
\numerone{3}
\section{Una tabella}
Le righe sottili sono la cosa che si rovina per prima quando un'immagine viene
ridimensionata male.

\vspace{1em}
\begin{center}
\begin{tabular}{|l|r|r|}
\hline
\textbf{Cosa} & \textbf{Prima} & \textbf{Dopo} \\
\hline
Scrollback terminali & 0 righe & 2000 righe \\
Scrollbar & assenti & dentro il frame \\
Comandi di esecuzione & globali & anche per progetto \\
Immagini & buffer vuoto & pixel veri \\
PDF & buffer vuoto & pagine sfogliabili \\
\hline
\end{tabular}
\end{center}
\newpage

% ---- 4 -------------------------------------------------------------------
\numerone{4}
\section{Una figura}
\begin{center}
\begin{tikzpicture}[scale=1.1]
  \draw[thick, fill=blue!10] (0,0) rectangle (8,4.5);
  \draw[thick, fill=green!10] (0,3.6) rectangle (8,4.5);
  \node at (4,4.05) {\small barra dei menu};
  \draw[thick, fill=yellow!10] (0,0) rectangle (2.2,3.6);
  \node[rotate=90] at (1.1,1.8) {\small albero file};
  \draw[thick] (2.2,1.6) -- (8,1.6);
  \node at (5.1,2.6) {\small editor};
  \node at (5.1,0.8) {\small terminali};
\end{tikzpicture}
\end{center}
Linee diagonali e curve sono il caso peggiore per i mezzi blocchi, e il caso in
cui i pixel veri si notano di più.
\newpage

% ---- 5 -------------------------------------------------------------------
\numerone{5}
\section{Il ciclo che conta}
Prova questo, che è il motivo per cui esiste tutto il resto:
\begin{enumerate}
  \item dividi l'editor con \texttt{Ctrl+L};
  \item tieni questo \texttt{.tex} da una parte e il \texttt{.pdf} dall'altra;
  \item cambia una parola qui sotto e salva con \texttt{Ctrl+S};
  \item premi \texttt{Ctrl+Shift+R}.
\end{enumerate}
La pagina accanto deve diventare quella appena composta, restando sulla pagina
che stavi guardando, senza chiudere e riaprire niente.

\vspace{1em}
\noindent\fbox{\parbox{0.9\textwidth}{\centering\large
Cambia questa frase e ricompila: dovresti vederla cambiare accanto.}}
\newpage

% ---- 6 -------------------------------------------------------------------
\numerone{6}
\section{Ultima pagina}
Da qui in avanti le frecce non devono fare niente: non un errore, non una pagina
bianca, proprio niente. Il conteggio dice che siamo alla fine, e anche se non lo
sapesse il rasterizzatore non produrrebbe nessun file per la pagina 7, che è il
modo in cui la fine si annuncia da sola.

\vspace{2em}
\noindent Torna indietro con $\leftarrow$ fino alla prima: anche lì le frecce
devono fermarsi senza lamentarsi.

\end{document}
